% !PRINTED 7
\weberpage{7}

\begin{center}
\S.\,2.

\gesperrt{The Divisors of Finite Groups.}
\end{center}

It is, as we have already seen in the first volume, a particularly important question whether a part of the elements of a group is itself again a group. For this it is necessary and sufficient that any two elements of this part yield under composition an element belonging again to the same part. This smaller group is then called a \gesperrt{divisor} or \gesperrt{factor}\footnote{Also called a subgroup.} of the whole group. The unit element ``1'' forms in every group a group by itself, and is therefore a divisor in every group.

We now restrict ourselves, as almost always in what follows, to \gesperrt{finite groups}, and shall assume that
\[
P = 1, a_1, a_2 \ldots a_{n-1}
\]
is a group of degree $n$ with the elements $a_0 = 1$, $a_1$, $a_2 \ldots a_{n-1}$, and that
\[
Q = 1, a_1, a_2 \ldots a_{\nu-1}
\]
is a divisor of $P$ of degree $\nu$. Every such divisor $Q$ must have the fundamental properties of a group, and therefore certainly contains the unit element ``1'', and with each of its elements $a_1$ at the same time the opposite $a_1^{-1}$. If $\nu < n$, we call $Q$ a \gesperrt{proper divisor} of $P$. If then $b$ is an element of $P$ not contained in $Q$, the elements
\[
Q_1 = b, a_1 b, a_2 b \ldots a_{\nu-1} b
\]
are all different from one another (\S.\,1, 3.) and all not contained in $Q$; for if, say, $a_1 b$ were contained in $Q$, then, since $Q$ is a group, $a_1^{-1} a_1 b = b$ would also have to be contained in $Q$, contrary to the hypothesis. If the whole group $P$ is not yet exhausted by $Q$ and $Q_1$, we take one of the remaining elements, which we may denote by $c$, and form
\[
Q_2 = c, a_1 c, a_2 c \ldots a_{\nu-1} c,
\]
and easily convince ourselves that the elements of $Q_2$ are all different not only from one another, but also from the elements of $Q$ and
%CONT>>
